% ------------------------------------------------------------------
% Draft section for Paper 6: Computational Complexity
% To be integrated after the experiment is run
% ------------------------------------------------------------------

\section{Computational Complexity of Structure}
\label{sec:complexity}

The channel-organized architecture is not only a modeling choice.
It is a compression with quantifiable computational consequences.

\subsection{The cost of the naive approach}

A transformer operating on $n$ genes uses self-attention with
$O(n^2)$ complexity per layer. For the full MSK-IMPACT panel
($n = 509$ genes), each attention layer computes $509^2 = 259{,}081$
pairwise interactions per patient. With $L$ layers and $H$ attention
heads, the total attention cost per patient is $O(L \cdot H \cdot n^2)$.

For $N = 43{,}872$ patients in the training set, the total
attention budget per epoch is:

\begin{equation}
\text{Attention ops (naive)} = N \cdot L \cdot H \cdot n^2
= 43{,}872 \cdot L \cdot H \cdot 259{,}081
\end{equation}

No published model has attempted this on MSK-IMPACT somatic
mutation data. Gene-level transformers exist for expression data
(GexBERT, Geneformer \cite{Geneformer2025}), where the input
modality is continuous-valued and information-dense.
Somatic mutation profiles are sparse and binary: most patients
carry mutations in fewer than 20 of 509 genes. Running
$O(n^2)$ attention over a vector that is 96\% zero is
computationally wasteful.

\subsection{The cost of the structured approach}

The channel-organized architecture compresses 509 genes into
$k = 16$ nodes (8 channels $\times$ 2 positions). Self-attention
operates over $k$ nodes: $16^2 = 256$ pairwise interactions per
layer per patient.

\begin{equation}
\text{Attention ops (structured)} = N \cdot L \cdot H \cdot k^2
= 43{,}872 \cdot L \cdot H \cdot 256
\end{equation}

The ratio:

\begin{equation}
\frac{n^2}{k^2} = \frac{509^2}{16^2} = \frac{259{,}081}{256}
\approx 1{,}012
\end{equation}

\textbf{The structured approach is $\sim$1{,}000$\times$ cheaper
in attention computation.} Same genes. Same data. Same patients.
The only difference is that 509 genes were organized into 16
channel-position nodes before the model saw them.

\subsection{Why this compression works}

The 1{,}000$\times$ reduction is not a lossy approximation.
It is a theory-derived factorization.

Papers~1--3 of this series prove that knowledge-maintaining systems
under bounded context organize into hierarchical coupling
structures. The channel map assigns each gene to its coupling
channel based on pathway membership. This is the graph structure
that Paper~1 proves is necessary for information preservation.

The naive transformer must \emph{discover} which gene--gene
interactions matter from the survival signal alone. It has
$n^2 = 259{,}081$ possible interactions to evaluate, of which
the vast majority are biologically meaningless (a mutation in
\emph{BRCA1} does not interact with a mutation in \emph{FOXA1}
through any known pathway). The model must learn to ignore
$>$99\% of its attention budget.

The structured approach provides the answer: genes within a
channel interact. Genes across channels interact at the channel
level. The attention budget is spent entirely on interactions
the theory predicts are meaningful.

\begin{table}[ht]
\centering
\caption{Computational comparison: naive gene-level transformer
vs.\ channel-organized transformer on MSK-IMPACT (509 genes,
43{,}872 patients).}
\label{tab:complexity}
\begin{tabular}{@{}lcc@{}}
\toprule
& \textbf{Naive} & \textbf{Structured} \\
\midrule
Input tokens per patient & 509 & 16 \\
Attention pairs per layer & 259{,}081 & 256 \\
Ratio & \multicolumn{2}{c}{$\sim$1{,}000$\times$} \\
\midrule
Parameters (embedding) & $509 \times d$ & $16 \times d$ \\
Parameters (attention) & $O(d^2)$ & $O(d^2)$ \\
\midrule
Information source & Data-driven & Theory-derived \\
Training data required & More & Less \\
Sparsity of input & $\sim$96\% zero & Dense (aggregated) \\
\bottomrule
\end{tabular}
\end{table}

\subsection{The hierarchical extension}

The channel-organized architecture admits a natural hierarchical
extension that further exploits the coupling structure:

\begin{enumerate}
\item \textbf{Level 1: Within-channel attention.} For each
channel, a small transformer operates over the genes in that
channel (e.g., 21 DDR genes, 14 Cell Cycle genes). This captures
within-channel interactions (e.g., BRCA1--BRCA2 co-mutation
patterns) at $O(m_c^2)$ cost where $m_c$ is the channel size.
Maximum channel size is 174 (PI3K/Growth in the expanded map).

\item \textbf{Level 2: Cross-channel attention.} The 8 (or 16)
channel-level embeddings from Level~1 are passed to a
cross-channel transformer. This captures which channel
combinations predict outcomes, at $O(k^2)$ cost.
\end{enumerate}

Total cost: $\sum_c m_c^2 + k^2$. For the expanded 509-gene,
8-channel map:

\begin{equation}
\sum_c m_c^2 + k^2 = 174^2 + 95^2 + 81^2 + 53^2 + 45^2
+ 38^2 + 15^2 + 8^2 + 16^2
= 51{,}270
\end{equation}

versus $509^2 = 259{,}081$ for the flat transformer.
A $5\times$ reduction even with full within-channel attention,
plus the architectural guarantee that cross-channel interactions
are learned at the right level of abstraction.

\subsection{Clinical deployment implication}

MSK-IMPACT sequences $>$50{,}000 patients. Targeted panels are
standard of care. A model that runs in $O(k^2)$ per patient can
be deployed as a real-time annotation on every sequencing report.
A model that runs in $O(n^2)$ cannot---not because the hardware
doesn't exist, but because the training data requirements and
overfitting risk scale with the parameter count.

The practical claim: theory-derived structure enables clinical
deployment that brute-force approaches do not. Not because the
brute-force approach is wrong, but because it is unnecessarily
expensive. The structure is free. It comes from the theory.
The compute savings come from using it.
